\section{Arithmetic}

\subsection{Topics}

\begin{itemize}
    \item Tropical semiring
    \item Tropical polynomial
    \item Newton polygon
\end{itemize}

\subsection{Definitions and results}

\begin{defin}[Tropical semiring]
    The \textit{min-plus tropical semiring} is
    \[\{\rr \cup \{\infty\},\oplus,\odot\},\]
    with tropical addition $a \oplus b = \min\{a,b\}$ and 
    tropical multiplication $a \odot b = a+b$.
\end{defin}

\begin{rmk}
    By the definition of a semiring, the following holds: 
    \begin{itemize}
        \item (Commutativity)
        \[x\oplus y=y\oplus x \quad \text{and}\quad x\odot y=y\odot x.\]
        \item (Distributive law)
        \[x\odot(y\oplus z)=x\odot y\oplus x\odot z.\]
        \item (Identity elements)
        \[x\oplus \infty=x\quad \text{and}\quad x\odot 0=0.\]
    \end{itemize}

    Note that the existence of an additive inverse is not guaranteed. For example, $6\oplus x=7$ has no solution for $x$.
\end{rmk}

\begin{defin}
    A \textit{tropical polynomial} is a finite linear combination of tropical monomials 
    \[p(x_1,\dots,x_n)=a\odot x_1^{i_1}x_2^{i_2}\cdots x_n^{i_n}\oplus b\odot x_1^{j_1}x_2^{j_2}\cdots x_n\oplus \cdots\]
    where the coefficients $a,b,\dots \in \rr$ and the exponents $i_1,j_1,\dots \in \zz$. 

    A tropical polynomial $p$ can be regarded as a function $p:\rr^n \to \rr$ by
    \[p(x_1,\dots,x_n)=\min(a+i_1x_1+\dots+i_nx_n,b+j_1x_1+\dots+j_nx_n,\dots)\]
\end{defin}

\begin{thm}
    The tropical polynomials in $n$ variables $x_1,\dots,x_n$ are precisely the piecewise-linear concave functions on $\rr^n$ with integer coefficients.
\end{thm}



\subsection{Examples}


\subsection{Exercises}
